\documentclass[11pt,a4paper]{article}
\usepackage[utf8]{inputenc}
\usepackage[T1]{fontenc}
\usepackage[spanish,es-noquoting]{babel}
\usepackage[margin=2.2cm]{geometry}
\usepackage{booktabs}
\usepackage{enumitem}
\usepackage{parskip}
\usepackage{amsmath}
\setlist{nosep,leftmargin=1.4em}
\pagestyle{empty}
\renewcommand{\arraystretch}{1.15}

\begin{document}

{\large\textbf{Plan de trabajo --- próximas etapas}}\\[2pt]
\textbf{Detección de familias de ransomware en base a archivos encriptados y notas de rescate}\\[2pt]
Romina Alfonzo, Carlos Urdapilleta \quad$\vert$\quad Tutor: Prof. Cristian Cappo \quad$\vert$\quad 12 de agosto de 2026

\hrulefill

\textbf{Punto de partida.} Los dos frentes cuentan con resultados consolidados y el capítulo de
resultados está redactado. La identificación de familia a partir de archivos cifrados alcanza
una exactitud de 0,910 sobre 29 familias, validada frente a tipos de documento no vistos
(0,879). La identificación a partir de notas de rescate alcanza un macro-F1 de 0,760 en el
escenario comparable con las herramientas existentes y de 0,435 ante plantillas nunca vistas.

El plan que sigue se ordena por dependencias, no por preferencia: la Etapa 1 condiciona los
resultados que se reportarán, y la Etapa 3 requiere que las anteriores estén cerradas.

\textbf{Fundamento del orden propuesto.} Tres verificaciones independientes indican que el
factor limitante ya no es metodológico:

\begin{center}
\small
\begin{tabular}{@{}llc@{}}
\toprule
\textbf{Verificación} & \textbf{Alcance} & \textbf{Efecto} \\
\midrule
Hiperparámetros del clasificador de notas & 840 configuraciones, búsqueda anidada & sin mejora \\
Abstracción de marcadores variables en las notas & 130 de 144 notas modificadas & sin mejora \\
Hiperparámetros de las características estadísticas & 4 combinaciones, hasta 2,5 h cada una & sin variación \\
\bottomrule
\end{tabular}
\end{center}

En consecuencia, el esfuerzo se concentra en la composición del corpus y no en la
optimización de los modelos.

\section*{Etapa 1 --- Ampliación del corpus de notas de rescate}

Es la única acción con capacidad de mejorar sustancialmente el resultado del segundo
protocolo de evaluación, porque ataca su causa.

El análisis de similitud mostró que las 146 notas del corpus corresponden a solo 95
contenidos distintos: cada familia reutiliza un repertorio reducido de plantillas. Siete
familias están representadas por una única plantilla, de modo que ningún método supervisado
puede generalizar en ellas.

\begin{itemize}
\item \textbf{Criterio propuesto:} al menos \textbf{tres plantillas distintas por familia},
en lugar de un número fijo de notas. La métrica se deriva del hallazgo anterior: incorporar
notas que repitan una plantilla ya presente no aporta capacidad de generalización.
\item \textbf{Fuentes identificadas:} se dispone de las direcciones concretas para catorce
familias con menos de cuatro notas, provenientes de portales de seguridad con
documentación verificable.
\item \textbf{Tarea complementaria:} auditar la procedencia de 37 notas cuya fuente quedó
registrada de forma ambigua, para que la sección de procedencia del capítulo de metodología
sea verificable en su totalidad.
\item \textbf{Duración:} es la etapa de duración menos previsible, por depender de la
disponibilidad pública de las notas. Se propone informar el avance por familia.
\end{itemize}

\section*{Etapa 2 --- Verificaciones pendientes}

Tareas acotadas, con costo computacional conocido, que cierran objeciones previsibles.

\begin{itemize}
\item \textbf{Naturaleza del bloque final en SUNCRYPT y NOTPETYA.} Ambas familias añaden
datos no aleatorios al final del archivo (entropía de cola de 4,78 y 6,58 frente a 7,44 en el
resto) y, sin embargo, no se identifican correctamente. La hipótesis es que ese bloque varía
en cada archivo ---una clave, un identificador de víctima o el tamaño original--- de modo que
constituye estructura pero no firma de familia. Se verificará comparando el contenido
hexadecimal de la cola en varias muestras. De confirmarse, permite redactar el diagnóstico
como hecho y no como conjetura.
\item \textbf{Incorporación del nombre de archivo como característica.} Lemmou et al. (2021)
obtienen un F de 0,920 clasificando únicamente por el nombre de la nota. En este trabajo se
descartó porque aproximadamente la mitad de los nombres del corpus fueron asignados por los
repositorios de origen y no por el atacante. Se propone anotar la procedencia de cada nombre
en el manifiesto y emplear como característica solo los auténticos.
\item \textbf{Reejecución sobre la base final.} Una vez ampliado el corpus y restauradas dos
notas de DHARMA que resultaron afectadas por el antivirus del equipo de trabajo, se repetirá
la corrida completa para que todas las cifras del documento provengan de un mismo conjunto.
\end{itemize}

\section*{Etapa 3 --- Redacción y cierre}

\begin{itemize}
\item Incorporar al capítulo de resultados los ocho conjuntos de mediciones ya obtenidas y
todavía no redactadas, entre ellas la validación de generalización a tipos de documento, la
localización de la información dentro del archivo y el diagnóstico de las familias de bajo
rendimiento.
\item Reformular en el marco teórico la delimitación del aporte respecto de Lemmou et al.
(2021), quienes identifican la familia mediante reglas y marcadores curados. La contribución
de este trabajo se define con precisión como el primer clasificador supervisado multiclase
sobre el contenido completo, con medición explícita de la generalización a variantes no vistas.
\item Completar la bibliografía con nueve referencias identificadas, entre ellas la del propio
conjunto de datos NapierOne, el trabajo sobre comparación de métodos de entropía del mismo
grupo, y el artículo de \textit{Electronics} sugerido en la reunión del 13/03/2024.
\item Redactar las conclusiones, el resumen y el \textit{abstract}, y completar el material
preliminar del documento. Incluir el agradecimiento al clúster computacional del NIDTEC, cuyo
reglamento de uso lo requiere.
\end{itemize}

\section*{Decisiones que requieren su criterio}

\begin{enumerate}[label=\alph*)]
\item \textbf{Procedencia de las notas.} Seis familias (BadRabbit, Chimera, Jigsaw, NotPetya,
WannaCry y CryptoLocker) no cuentan con archivo bruto público disponible y se cubrieron con
transcripciones documentadas de portales de seguridad. ¿Se consideran fuentes admisibles,
declarándolo como limitación de procedencia?
\item \textbf{Criterio de ampliación.} ¿Resulta adecuado el criterio de tres plantillas
distintas por familia en lugar de un número fijo de notas?
\item \textbf{Estructura del documento.} ¿Conviene reorganizar los capítulos de modo que cada
objetivo específico corresponda a una sección, según lo sugerido el 08/05/2024?
\item \textbf{Alcance.} El clúster dispone de unidades de procesamiento gráfico. ¿Corresponde
incorporar una comparación con modelos basados en \textit{transformers} como sección
adicional, o dejarla planteada como trabajo futuro, considerando que el corpus consta de 146
documentos y tales modelos tienden al sobreajuste en conjuntos de ese tamaño?
\item \textbf{Familias con una sola plantilla.} Para las siete familias en esa situación,
ningún método supervisado clásico puede generalizar. ¿Incorporar el enfoque de aprendizaje
\textit{few-shot} del repositorio que nos compartió el 06/06/2024, o dejarlo como trabajo
futuro?
\end{enumerate}

\section*{Sobre las seis familias que el análisis de archivos no resuelve}

Conviene señalar que las seis familias cuya identificación por archivos cifrados resulta
deficiente ---SUNCRYPT, WASTEDLOCKER, CRYPTOLOCKER, DARKSIDE, JIGSAW y NOTPETYA--- refuerzan
el planteo de investigar ambos artefactos por separado: son precisamente los casos en que la
nota de rescate constituye la única vía de identificación disponible. La independencia de los
dos frentes, mantenida a lo largo del trabajo, resulta así justificada por los datos y no
solo por el diseño.

\end{document}
